\documentclass[11pt]{article}
\usepackage{fontspec}
\usepackage[french]{babel}
\usepackage{amsmath}
\usepackage{amssymb}
\usepackage{amsthm}
\usepackage{mathrsfs}
\usepackage{tikz-cd}
\usepackage{enumitem}
\usetikzlibrary{babel}
\tikzcdset{every diagram/.append style={baseline={([yshift=-.5ex]current bounding box.center)}}}

\newcommand{\printedpage}[1]{\marginpar{\footnotesize p.~#1}}
\newcommand{\msp}[1]{\marginpar{\footnotesize [#1]}}
\newcommand{\Hom}{\mathrm{Hom}}
\newcommand{\id}{\mathrm{id}}
\newcommand{\Winf}{\mathbf{W}_{\infty}}

\begin{document}

\printedpage{5}
% head: A. Grothendieck : Dérivateurs, Ch. XII, Édition des Universités de Montpellier II et Paris VII
% foot: Version 20 juin 2003

dans $U'$, et les hypothèses a), b) de la page 1 assurent que la somme amalgamée $\varinjlim$ de ce diagramme dans $U'^{\wedge}$ est $e_{U'}$, i.e. que
\[ p_!(e_U) = e_{U'}. \]
Il s'ensuit que dans la somme amalgamée, $X'$ envisagée dans $(*)$, l'unique élément $t'$ de $T'$ est objet final de $X'$. On a vu que $U'$ est 1-connexe, mais non 2-connexe si le groupe abélien enveloppant de $G$ (i.e. $G/$(groupe des commutateurs) dans le cas où $G$ est un groupe) n'est pas réduit à 1. L'exemple le plus économique s'obtient en prenant $G = \{\pm 1\}$ le groupe à deux éléments. C'est donc un cas où dans $(*)$ on a
\[ i \in W \ \ \text{(et même $i$ asphérique),} \qquad \text{mais } i' \notin W, \]
et qui a l'avantage sur l'exemple plus sophistiqué de la page 103 (prop. 10) de ne mettre en jeu que des catégories finies -- et même fort petites !

Procédons maintenant comme dans la proposition 11, en désignant par $U''$ la catégorie déduite de $U'$ en lui ajoutant un objet
\msp{page 7}
final, et soit
\[ p' : U' \longrightarrow U'' \]
l'inclusion, de sorte que
\[
\begin{aligned}
p'_!(e_{U'}) \quad &= \quad (p'p)_!(e_U) = U' \quad && \text{en tant qu'objet dans } U''^{\wedge}\\
& && \text{(sous-préfaisceau du préfaisceau final } e_{U''}\text{).}
\end{aligned}
\]
Considérons
\[
\begin{tikzcd}
U \arrow[r, hook, "i"] \arrow[d, "{p''=p'p}"'] \arrow[dr, phantom, "\begin{smallmatrix}\text{cocart.}\\ \text{et cart.}\end{smallmatrix}"] & X \arrow[d, "q''"'] & T \arrow[l, hook', dotted, "j"'] \arrow[d, dotted, "\text{iso}"] \\
U'' \arrow[r, hook, "i''"] & X'' & T''. \arrow[l, hook', dotted, "j''"']
\end{tikzcd}
\]
Cette fois $U$, $X$ \textit{et} $U''$ sont $\Winf$-asphériques (et même contractiles toutes trois), donc $i$, $p''$ sont dans $W$, pourtant $X''$ n'est pas $\Winf$-asphérique (donc $i'', q'' \notin \Winf$), car pour qu'elle le soit, i.e. pour que $i'' \in \Winf$, il faudrait (prop. 7, cor. page 85 (\footnote[4]{pas encore démontré !})) que $i''$ soit $\Winf$-asphérique, i.e. que $U''/t''$ soit $\Winf$-asphérique. Mais par construction $U''/t'' \simeq U'$, qui n'est \textit{pas} $\Winf$-asphérique.

Voilà un exemple encore plus simple d'un foncteur $f : \Psi \longrightarrow X$ tel que $f_!(e_{\Psi})$ ne soit pas $\Winf$-asphérique, i.e. d'un diagramme
\[
\begin{tikzcd}
& b \\
a \arrow[ur, "\beta_0"] \arrow[dr, "\gamma_0"'] & \\
& c
\end{tikzcd}
\]
\printedpage{6}
% head: A. Grothendieck : Dérivateurs, Ch. XII, Édition des Universités de Montpellier II et Paris VII
% foot: Version 20 juin 2003
\msp{page 8}
dans une catégorie $V$, telle que la limite inductive $L = b \amalg_a c$ \textit{dans} $V^{\wedge}$ soit un objet non asphérique de $V^{\wedge}$, i.e. $V/L$ soit non asphérique. Je sais qu'il faut, pour un tel contre-exemple, que $a \longrightarrow b$ et $a \longrightarrow c$ ne soient pas des monomorphismes. L'idée, c'est de prendre un exemple universel où ces flèches ne soient pas \texttt{[des]} monomorphismes, en prenant
\msp{page 9}
\[
V = \left(
\begin{tikzcd}
& & b \\
0 \arrow[r, shift left, "\lambda"] \arrow[r, shift right, "\mu"'] & a \arrow[ur, "\beta"] \arrow[dr, "\gamma"'] & \\
& & c
\end{tikzcd}
\right)
\]
avec
\[ \beta\lambda = \beta\mu \qquad \text{et} \qquad \gamma\lambda = \gamma\mu \]
les deux seules flèches non identiques, en plus des quatre flèches marquées $\beta$, $\gamma$, $\lambda$, $\mu$.

On trouve que
\[ b \amalg_a c \xrightarrow{\ \sim\ } e_X \qquad \text{(somme amalgamée dans } V^{\wedge}\text{).} \] % [sic] e_X -> e_V, the category considered here is V

Pourtant $V$ n'est pas asphérique, ni même 1-connexe : il est homotope à la sous-catégorie ouverte $V_0 = (0 \overset{\lambda}{\underset{\mu}{\rightrightarrows}} a)$, car on a une rétraction $V \xrightarrow{\ r\ } V_0$ (envoyant $b$ et $c$ sur $a$), et une flèche $ir \longrightarrow \id_V$ (où $i : X_0 \hookrightarrow X$ est l'inclusion). Donc $V$ a le même type d'homotopie \textit{que} $V_0$, qui est un cercle. Donc on a intérêt à prendre ce $V$, où $\tilde{V}$ déduit de $V$ en rajoutant un élément final $e$, % [sic] i : X_0 \hookrightarrow X -> i : V_0 \hookrightarrow V, the inclusion of V_0 in V is meant
\[
\tilde{V} = \left(
\begin{tikzcd}
& & b \arrow[dr] & \\
0 \arrow[r, shift left] \arrow[r, shift right] & a \arrow[ur] \arrow[dr] & & e \\
& & c \arrow[ur] &
\end{tikzcd}
\right),
\]
comme $U'$, pour les contre-exemples pages 6 et 7.

\msp{page 10}
\subsection*{2.6. Contre-exemples pour la totale asphéricité.}

Je me suis finalement convaincu que
\[
\Psi = \left\{
\begin{tikzcd}
& b_0 \\
a_0 \arrow[ur] \arrow[dr] & \\
& c_0
\end{tikzcd}
\right\}
\]
\printedpage{7}
% head: A. Grothendieck : Dérivateurs, Ch. XII, Édition des Universités de Montpellier II et Paris VII
% foot: Version 20 juin 2003
n'est \textit{pas} totalement asphérique (pour le localiseur $\Winf$ -- dans le présent numéro, on supposera $W = \Winf$), ou ce qui revient au même, qu'on peut avoir un diagramme de type $\Psi$
\[
\begin{tikzcd}
& b \\
a \arrow[ur, "\beta"] \arrow[dr, "\gamma"'] & \\
& c
\end{tikzcd}
\tag*{$(*)$}
\]
dans une catégorie $Z \in \mathrm{Ob}\,\mathrm{Cat}$, dont la limite inductive $b \amalg_a c$ dans $Z^{\wedge}$ soit le préfaisceau final $e_Z$, i.e. tel que l'on ait
\[ b(z) \amalg_{a(z)} c(z) \xrightarrow{\ \sim\ } \text{ensemble ponctuel} \]
pour tout objet $z$ de $Z$ (où $a(z) = \Hom_Z(a,z)$ \texttt{[plutôt} $\Hom_Z(z,a)$\texttt{]}) sans que pour autant $Z$ soit asphérique, ni même 1-connexe (\footnote[5]{faux -- il est 1-connexe forcément (\footnotemark[6]), mais non 2-connexe.}).\footnotetext[6]{cf. commentaire XII, p. 101.} L'exemple le plus simple que j'ai trouvé, est l'exemple ‘universel’ d'un diagramme $(*)$ où ni l'une ni l'autre des deux flèches $a \longrightarrow b$, $a \longrightarrow c$ n'est \texttt{[un]} monomorphisme (car si elle l'est, alors on peut prouver que $Z$ est bel et bien asphérique, comme nous verrons par la suite). C'est la catégorie
\[
Z = \left\{
\begin{tikzcd}
& & b_0 \\
0_0 \arrow[r, shift left, "\lambda_0"] \arrow[r, shift right, "\mu_0"'] & a_0 \arrow[ur, "\beta_0"] \arrow[dr, "\gamma_0"'] & \\
& & c_0
\end{tikzcd}
\right\}
\]
\msp{page 11}
déduite de $\Psi$ en lui rajoutant un objet $0_0$, relié à $a_0$ par deux flèches distinctes $\lambda_0$, $\mu_0$, telles que
\[ \beta_0\lambda_0 = \beta_0\mu_0 \qquad \text{et} \qquad \gamma_0\lambda_0 = \gamma_0\mu_0, \]
ces deux flèches composées étant, avec les quatre flèches $\beta_0$, $\gamma_0$, $\lambda_0$, $\mu_0$ marquées sur le diagramme, les seules (six) flèches non identiques de $Z$. $\Psi$ immergée dans $Z$ comme un fermé, de complémentaire la catégorie ponctuelle réduite à $0_0$. (NB L'immersion fermée $\Psi \longrightarrow Z$ n'est pas immersion de Dwyer, i.e. (comme $Z$ est le plus petit ouvert de $Z$ contenant $\Psi$) n'est pas coasphérique, car la catégorie $0\backslash\Psi$, cofibrée sur $\Psi$ avec fibres discrètes, de cardinal 2 en $a_0$ et 1 en $b_0$ et $c_0$, est la catégorie
\[
\begin{tikzcd}
& \tilde{b}_0 \\
a'_0 \arrow[ur] \arrow[ddr] & \\
a''_0 \arrow[uur, crossing over] \arrow[dr] & \\
& \tilde{c}_0
\end{tikzcd}
\]
\printedpage{8}
% head: A. Grothendieck : Dérivateurs, Ch. XII, Édition des Universités de Montpellier II et Paris VII
% foot: Version 20 juin 2003
i.e.
\[
\begin{tikzcd}
& {} & \\
{} \arrow[ur] \arrow[dr] & & {} \arrow[ul] \arrow[dl] \\
& {} &
\end{tikzcd}
\quad ,
\]
qui a le type d'homotopie du cercle.) Ni l'immersion ouverte $\{0_0\} \longrightarrow Z$, car $\{0_0\}/a_0$ est la catégorie discrète de cardinalité 2.

L'inclusion $i : Z_0 \longrightarrow Z$ de
\[ Z_0 \ = \ (\ 0_0 \overset{\lambda_0}{\underset{\mu_0}{\rightrightarrows}} a_0\ ) \]
dans $Z$ est un homotopisme (\footnote[7]{faux, il n'y a pas de rétraction du tout! En fait, $Z$ est la \textit{suspension} de la $\rightrightarrows$, qui a le type d'homotopie (topologique) du cercle, et il a donc le type d'homotopie de la 2-sphère.}) (car il existe une rétraction évidente $r$ de $Z$ sur $Z_0$ (\footnotemark[7]), et une flèche $ir \longrightarrow \id_Z$), donc $Z$ est homotope à $Z_0$, qui a le type d'homotopie du cercle, donc n'est pas 1-connexe.

\msp{page 12}
D'autre part, la vérification de
\[ b_0 \amalg_{a_0} c_0 \ = \ e_Z \qquad \text{dans } Z^{\wedge} \]
est immédiate.

On va utiliser cette flèche
\[ \Psi \longrightarrow Z \]
pour construire des exemples intéressants (et déchirants) de sommes amalgamées
\[
\begin{tikzcd}
U \arrow[r, hook, "i"] \arrow[d, "p"'] & X \arrow[d, "q"'] \\
U' \arrow[r, hook, "i'"] & X',
\end{tikzcd}
\]
avec $i$ immersion ouverte, telle qu'on ait même
\[ i :\ U \longrightarrow \overline{U} \qquad \text{est asphérique, donc dans } W, \]
alors que pourtant on aura
\[ i' :\ U' \longrightarrow \overline{U'} \qquad \notin W, \]
et ce qui mieux (ou pis) est,
\[ p \in W \qquad \text{mais} \qquad q \notin W \]
\printedpage{9}
% head: A. Grothendieck : Dérivateurs, Ch. XII, Édition des Universités de Montpellier II et Paris VII
% foot: Version 20 juin 2003
(et éventuellement, en même temps, $i \in W$ mais $i' \notin W$). De plus (contrairement aux constructions des propositions 10, 11 (p. 101-107)), les catégories envisagées pourront toutes être finies. Je suis particulièrement intéressé par le cas où $X$ est un ‘carré’
\[ X = \Phi_n \times \Phi_m \qquad (n,m \geq 2), \]
et où $U$ est son ‘bord’
\[ U = \dot{X} \overset{\text{déf}}{=} (\dot{\Phi}_n \times \Phi_m) \cup (\Phi_n \times \dot{\Phi}_m). \]

\msp{page 13}
\[
\begin{tikzcd}
0 \arrow[r] \arrow[d] & a \arrow[d] & a' \arrow[l] \arrow[r] \arrow[d] & \cdots \\
b \arrow[r] & t_1 & \cdots \arrow[l] & \\
b' \arrow[u] \arrow[r] \arrow[d] & {} \arrow[u] & \cdots \arrow[l] \arrow[u] & \\
\cdots & & &
\end{tikzcd}
\]

\textbf{Proposition.} \textit{Soient $n,m \geq 2$, et considérons l'inclusion $U = \partial X \overset{i}{\hookrightarrow} X$ du bord du ‘carré’ $X = \Phi_n \times \Phi_m$ dans $X$. C'est une immersion ouverte $W$-asphérique, de façon précise, les $U/t$ ($t \in \dot{U} = \bar{U}\backslash U$) sont isomorphes à $e$ (catégorie ponctuelle) ou à $\Phi =$ \begin{tikzcd}[ampersand replacement=\&, row sep=small, column sep=small] \& {} \arrow[dl] \\ {} \& \\ \& {} \arrow[ul] \end{tikzcd} (catégorie avec élément final), sauf pour les quatre éléments $t_i$ des quatre ‘coins’ de $\dot{U}$, où c'est la catégorie}
\[
\tilde{\Psi} = \left(
\begin{tikzcd}
& & a' \arrow[dl] \\
& a & \\
0 \arrow[ur] \arrow[dr] & & \\
& b & \\
& & b' \arrow[ul]
\end{tikzcd}
\right)
\]
\textit{n'ayant pas d'élément final. Donc $U \xrightarrow{\ i\ } X$ n'est pas immersion ouverte de Dwyer (mais $U \longrightarrow X\backslash\{t_1,t_2,t_3,t_4\}$ l'est \ldots). Ceci posé, je dis que $i : U \longrightarrow X$ n'est pas une $W$-cofibration.}

\end{document}
